

% \begin{question}
% 	It is straightforward to see $D^B_{r L} = r D^B_L$, when $r \in \mathbb{Q}_{> 0}$. What is the relationship between $D^B_{L_1 + L_2}$ and $D^B_{L_1}, D^B_{L_2}$, when $L_1, L_2$ are two ample $\mathbb{Q}$-line bundles?
% \end{question}

We are going to show the $\mathbb{Q}$-divisor $D^B_L := \bc_{\dd \mu}(\Rls{L}^B)$ is the spherical analogy of $D^T_L$ in \Cref{sec:comp:tor}. Similar to \Cref{thm:toric}, our main result is
\begin{theorem}\label[theorem]{thm:sph}
	Let $X$ be a spherical $G$-variety, and let $L$ be a $\mathbb{Q}$-line bundle on $X$, the $G$-equivariant stability threshold of $L$ satisfies
	\begin{equation*}
		\delta_G(L) = \inf_{v \in \DVal^G_X} {A_X(v) \over v(D^B_L)}.
	\end{equation*}
\end{theorem}

Given a polarized spherical $G$-variety $(X, L)$ and a $G$-linearization $\alpha$, for every $m \in \mathbb{N}(L, \alpha)$ such that $m L$ is very ample, we can decompose the finite-dimensional $G$-module $H^0(X, m L)$ into
\begin{equation}\label{eq:decompSph}
	H^0(X, m L) = \bigoplus_{\lambda \in P_m(X, L, \alpha)} H^0(X, m L)_\lambda .
\end{equation}
For simplicity, let $V := H^0(X, m L)$ (thus $V_\lambda = H^0(X, m L)_\lambda$), and $N_{m, \lambda} := \dim(V_\lambda)$.

For every character $\lambda \in P_m(X, L, \alpha)$, we have a unique (up to scaling) non-zero section $e^B_\lambda \in V^{(B)}_\lambda$.

\begin{lemma}\label[lemma]{lem:GvalBYeb}
	Given two valuations $v, v' \in \DVal^G_X$, if for every character $\lambda \in P_m(X, L, \alpha)$, we have $v(e^B_\lambda) = v'(e^B_\lambda)$, then $v = v'$.
\end{lemma}
\begin{proof}
	It is a simple corollary of \Cref{lem:moneSPAN} and \Cref{def:valcone}. Note that we have assumed $m L$ to be very ample.
\end{proof}

\begin{lemma}\label[lemma]{lem:GvalCONST}
	Every valuation $v \in \DVal^G_X$ is constant on $V_\lambda \setminus \{0\}$ for every character $\lambda \in P_m(X, L, \alpha)$; more precisely, the equality $v(s) = v(e^B_\lambda)$ holds for every section $s \in V_\lambda \setminus \{0\}$.
\end{lemma}
\begin{proof}
	For every $v \in \DVal^G_X$, we can define a $G$-invariant valuation
	\begin{equation*}
		\hat{v}(s) := \min\limits_{s_\lambda \neq 0} v(e^B_\lambda)
	\end{equation*}
	where $s = \sum_{\lambda \in P_m(X, L, \alpha)} s_\lambda \in \bigoplus_{\lambda \in P_m(X, L, \alpha)} V_\lambda$. By \Cref{lem:GvalBYeb}, we have $v = \hat{v}$.

	Therefore, the equality $v(e^B_\lambda) = \hat{v}(e^B_\lambda) = \hat{v}(s) = v(s)$ holds for every nonzero section $s \in V_\lambda \setminus \{0\}$.
\end{proof}

\begin{lemma}\label[lemma]{lem:GvalCPTdecomp}
	The decomposition \Cref{eq:decompSph} is compatible (see \Cref{sec:comp}) with any valuation $v \in \DVal^G_X$.
\end{lemma}
\begin{proof}
	According to the proof of \Cref{lem:GvalCONST}, every $v \in \DVal^G_X$ satisfies
	\begin{equation*}
		v(s) = \min_{s_\lambda \neq 0} v(s_\lambda)
	\end{equation*}
	where $s = \sum s_\lambda$.
\end{proof}

We are now well-prepared to prove the main theorem.

\begin{proof}[Proof of \Cref{thm:sph}]
	For every character $\lambda \in P_m(X, L, \alpha)$, we pick a basis $\{s_{\lambda, j}\}_j \subseteq V_\lambda$. For every valuation $v \in \DVal^G_X$, by \Cref{lem:GvalCONST}, the valuation $v$ is constant on $V_\lambda$, so the basis $\{s_{\lambda, j}\}_j$ is compatible with $v$ on $V_\lambda$. Now $\{s_{\lambda, j}\}_{\lambda, j}$ is a basis of $V_m$ which is compatible with every $v \in \DVal^G_X$, and we denote the corresponding $m$-basis type $\mathbb{Q}$-divisor
	\begin{equation*}
		D_m = {1 \over N_m} \sum_{\lambda \in P_m(X, L, \alpha)} \sum_{j = 1}^{N_{m, \lambda}} \divQ(s_{\lambda, j}).
	\end{equation*}
	Consequently, by \Cref{lem:Smdef}, we have,
	\begin{equation*}
		\delta_{m, G}(L) = \inf_{v \in \GDValo_X} {A_X(v) \over v(D_m)}.
	\end{equation*}

	It is challenging to find the limit of $(D_m)_m$, as $\Qls{L}$ is infinite-dimensional. That is why we introduce the $B$-invariant $\mathbb{Q}$-divisor  $D^B_m := {1 \over N_m} \sum_\lambda N_{m, \lambda} \divQ(e^B_\lambda) = \bc_{\mu_m}(\Rls{L}^B)$. We have,
	\begin{equation*}
		v(D_m) = {1 \over m N_m} \sum_\lambda \sum_{j = 1}^{N_{m, \lambda}} v(s_{\lambda, j}) = {1 \over m N_m} \sum_\lambda N_{m, \lambda} \: v(e^B_\lambda) = v(D^B_m).
	\end{equation*}
	The sequence $(D^B_m)_m$ is convergent to $D^B_L = \bc_{\mu}(\Rls{L}^B)$ in the finite dimensional topological space $\Rls{L}^B$. As $v$ is linear (thus continuous) on $\Rls{L}^B$, we finally get,
	\begin{equation*}
		\delta_G(L) = \inf_{v \in \GDValo_X} {A_X(v) \over v(D^B_L)}.
	\end{equation*}
\end{proof}

\subsection{Comparison with Brion's \textit{B}-invariant anticanonical divisor}\label{sec:KstSph:Brion}


Similar to toric varieties, for a spherical $G$-variety $X$, there are only finitely many $B$-invariant prime divisors, see \cite[Theorem 2.2.2]{Pezzini_lec_sph_wond_vari}. Let
\begin{itemize}
    \item $X_1, \dots, X_k$ be the $G$-invariant prime divisors on $X$, and
    \item $Y_1, \dots, Y_{k'}$ be the $B$-invariant but not $G$-invariant prime divisors on $X$.
\end{itemize}
\begin{remark}
    The prime divisors $Y_1, \dots, Y_{k'}$ are called \emph{colors} of a spherical variety.
\end{remark}

As the $\mathbb{Q}$-divisor $D^B_X$ is $B$-invariant, we can find the coefficients of $D^B_X$ by 
\begin{equation}
    D^B_X = \sum_{i = 1}^{k'} \alpha_i Y_i + \sum_{j = 1}^{k} \beta_j X_j
\end{equation}
where $\alpha_i, \beta_j \in \mathbb{Q}_{\geq 0}$.

Similar to \Cref{cor:tor:Kss}, we have a weaker result for spherical varieties,
\begin{corollary}\label[corollary]{cor:sph:Kss}
	If a $\mathbb{Q}$-Fano spherical variety $X$ is K-semistable, then we have the inequality $\beta_j \leq 1$ for every $j = 1, \dots, k$.
\end{corollary}
\begin{proof}
    This follows directly from \Cref{thm:sph} and \Cref{thm:Zhuang}. For every $j$, the coefficient $\beta_j = \ord_{X_j}(D^B_X)$, and as divisors $X_1, \dots, X_k$ are prime divisors on $X$, we have $A_X(\ord_{X_j}) = 1$.
\end{proof}

There is another natural $B$-invariant representative for anticanonical $\mathbb{Q}$-divisors, which is indeed a $\mathbb{Q}$-Cartier divisor, for a (klt) Fano spherical $G$-variety $X$.
\begin{theorem}[c.f., {\cite[Theorem 1.2]{GH_gor_sph_fano_var}}, based on {\cite[\S4.1, 4.2]{Brion_cur_div_sph_var}}]
    There is a $B$-invariant anticanonical divisor of the form
    \begin{equation*}
        -\cano_X := \sum_{i = 1}^{k'} m_i Y_i + \sum_{j = 1}^{k} X_j,
    \end{equation*}
    where the coefficients $m_j \in \mathbb{Z}_{> 0}$ depend only on the open $G$-orbit of variety $X$.
\end{theorem}
\begin{remark*}
    We call this divisor \emph{Brion's anticanonical divisor}. For toric varieties, Brion's anticanonical divisor is the same as the standard anticanonical divisor $X_1 + \dots + X_k$.
\end{remark*}

Following \Cref{sec:prelim:bir:dis}, we can use Brion's anticanonical divisor to define the log discrepancy.

There is a deep relationship between Brion's anticanonical divisor and the $\mathbb{Q}$-divisor $D^B_X$ for K-semistable varieties.
\begin{corollary}
    Given a $\mathbb{Q}$-Fano spherical $G$-variety $X$, $X$ is K-semistable iff for every valuation $v \in \GDValo_X$, the equality $v(D^B_X) \leq v(-\cano_X)$ holds. Specifically, if variety $X$ is horospherical, which means the valuation cone is the whole space, then $X$ is K-semistable iff $D^B_X = -\cano_X$.
\end{corollary}
\begin{remark*}
    This is the reformulation of Delcroix's criteria (see \cite[Theorem 5.3 and Corollary 5.7]{Delcroix_Kst_Fano_spher}). Note that the anticanonical divisor $-\cano_X \in \Qls{L}^B$ here plays the same role as the point $2 \rho_Q \in \Delta^+$ in \textit{loc.~cit.}
\end{remark*}
\begin{proof}
    According to \cite[Proof of Proposition 5.2]{Pas_sing_sphe}, for every valuation $v \in \GDValo$, we have $v(-\cano_X) = A_X(v)$. The result then follows directly from \Cref{thm:sph}.
\end{proof}

Finally, let us study the coefficient $\alpha_i$ of the $\mathbb{Q}$-divisor $D^B_X$. If the image of valuation $\ord_{Y_i}$ is in the valuation cone, then we have the inequality $\alpha_i \leq m_i$. In general, coefficient $\alpha_i$ may be larger than $m_i$.
\begin{example}
    Let $G := \SL_2(\mathbb{C})$, and the Borel subgroup $B$ is the subgroup of upper triangular matrices. The variety $\mathbb{P}^1 \times \mathbb{P}^1$ is spherical under the diagonal action. The variety $X$ has three $B$-invariant prime divisors,
    \begin{equation*}
    \begin{aligned}
        Y_1 &:= \{[1: 0]\} \times \mathbb{P}^1, \qquad Y_2 := \mathbb{P}^1 \times \{[1: 0]\}, \\
        X_1 &:= \{(x, x) \in \mathbb{P}^1 \times \mathbb{P}^1 \mid x \in \mathbb{P}^1 \}.
    \end{aligned}
    \end{equation*}
    The anticanonical bundle is $\stsh_X(2, 2)$, so that $-\cano_X = Y_1 + Y_2 + X_1$, and according to \cite[Example 5.10]{Delcroix_Kst_Fano_spher}, we have $D^B_X = {4 \over 3} Y_1 + {4 \over 3} Y_2 + {2 \over 3} X_1$.
    
    \begin{figure*}[h]
        \begin{tikzpicture}[
    dot/.style={circle, fill, inner sep=1.5pt}, % 定义点的样式
    label distance=2pt
]
    % 2. 绘制主线段
    \coordinate (A) at (0,0);
    \coordinate (B) at (6,0);
    \draw[thick] (A) -- (B);

    % 3. 标注并绘制点
    % 端点
    \node[dot, label=above:{$2 X_1$}] at (A) {};
    \node[dot, label=above:{$2 Y_1 + 2 Y_2$}] at (B) {};

    % 中点 (1/2)
    \node[dot, label=above:{$-\cano_X$}] (M) at ($(A)!0.5!(B)$) {};

    % 2/3 点
    \node[dot, label=below:{$D^B_X$}] (T) at ($(A)!0.6667!(B)$) {};

\end{tikzpicture}
        \caption{Rational linear series $\Qls{{-\cano_X}}^B$.}\label[figure]{fig:P1xP1}
    \end{figure*}

    In this case, we have the inequality $\alpha_1 = {4 \over 3} > 1 = m_1$.
\end{example}


